
\documentclass[12pt]{article}
\usepackage[a4paper, total={6.5in, 9.5in}]{geometry}
\usepackage{graphicx}			
\usepackage{amssymb}
\usepackage{amsmath}			
\usepackage{fancyhdr}
\usepackage{enumerate}       
\usepackage{dirtytalk} 
\usepackage[english]{babel}
\usepackage{amsthm}
\usepackage{xcolor}
\usepackage[shortlabels]{enumitem}
\setlength{\parindent}{0pt}
\setlength{\parskip}{0.8em}

\newtheorem{theorem}{Theorem}[section]
\newtheorem{corollary}{Corollary}[theorem]
\newtheorem{lemma}[theorem]{Lemma}
\newtheorem{definition}[theorem]{Definition}


\begin{document}
% \renewcommand{\thesection}{Exercise \arabic{section}}

\setcounter{section}{0} % Sets the section counter to 4, so the next section is 5
\renewcommand{\thesection}{\arabic{section}}


\parindent = 0pt


      {\bf Analysis 1 HW 3} 


%==================


\section{}
 {\bf Definitions}\\
$(x_{n})$ is a sequence of real numbers satisfying $3x_{n+1} = {x_{n}^{3}-2}$ for all $n\in \mathbb{N}$.

      {\bf Solution}\\
If $(x_{n})\to l$ for some $l$, then
\begin{align*}
               & \lim_{ n \to \infty } (3x_{n+1}) = \lim_{ n \to \infty } (x_{n}^{3}-2) \\
      \implies & 3l=l^{3}-2
\end{align*}
courtesy the algebraic limit theorem. This yields $l\in \{ 2, -1 \}$.

Now, consider the successive differences:
\begin{align}
      x_{n+1}-x_{n} & = \frac{{x_{n}^{3}-3x_{n}-2}}{3}
\end{align}

\begin{enumerate}[(i)]
      \item If $x_{i} > 2$, $x_{i}^{3}-3x_{i}-2 = x_{i}(x_{i}^{2}-3)-2>2(4-3)-2=0$. So, $x_{i+1}-x_{i}>0$. Therefore, if $x_1>2$, the sequence is increasing and cannot converge to 2 or -1, so it must diverge.
      \item If $x_{i} = 2$, $x_{i+1} -x_{i}=0$. Thus, if $x_1 =2$ the sequence is constant at 2, and it converges to 2.
      \item If $-1<x_{i}<2$, we make the following claims:
            \begin{enumerate}
                  \item {\bf Claim: }If $-1<x_{i}<2$, then $x_{i+1}-x_{i}\leq0$.\\
                        It suffices to show that $y^{3}-3y-2\leq0$ for all $y<2$.
                        Consider the expression $t(6-t)-9$, $t>0$. If $t>6$, it is clear that $t(6-t)-9 <0$. If $0<t\le 6$ , $t$ and $6-t$ are non negative. From the AM-GM inequality, we have $9\le t(6-t) \implies$  $t(6-t)-9\le 0$. Thus, $t(t(6-t)-9)\le 0$ is true for all $t>0$. Now, let $y=2-t$. Plugging $t=2-y$ into the above inequality, we get ${{{y}^{3}}-{3y}-2}\le 0$ for all $y<2$.
                  \item {\bf Claim: }If $-1<x_{i}<2$, then $-1<x_{i+1}<2$.\\
                        \begin{equation*}
                              x_{i}<2
                              \implies x_{i}^{3}<8
                              \implies  x_{i}^{3}-2<6
                              \implies  \frac{{x_{i}^{3}-2}}{3}=x_{i+1}<2
                        \end{equation*}

                        \begin{equation*}
                              -1<x_{i}
                              \implies  -1<x_{i}^{3}
                              \implies  -3<x_{i}^{3}-2
                              \implies  -1< \frac{{x_{i}^{3}-2}}{3} = x_{i+1}
                        \end{equation*}
            \end{enumerate}
            Thus, if $x_{1}$ lies between $-1$ and $2$, then $(x_{n})$ is bounded and monotone decreasing. By the monotone convergence theorem, the sequence must converge, and since it cannot converge to 2, it must converge to -1.
      \item If $x_{i} = -1$, $x_{i+1} -x_{i}=0$. Thus, if $x_1 = -1$ the resulting sequence is constant at -1, and it converges to -1.
      \item For $x_i < -1$, we have $x_i^3 - 3x_i - 2 \le 0$ from (iii)(a) (while claimed for $-1<x_i<2$, we proved it for all $x_i<2$). Thus, for $x_1<-1$ the sequence is decreasing. Since the sequence cannot converge to either -1 or 2, it must diverge to negative infinity.
\end{enumerate}

\section{}
 {\bf Definitions}\\
$\{ (a_{i}, b_{i}) \}$ is a sequence in $\mathbb{R}^{2}$ with $a_{i}\le b_{i}$ such that for each $i$, the closed interval $[a_{i}, b_{i}]$ contains the closed interval $[a_{i+1}, b_{i+1}]$.\\
$\sup a_{n} = \sup \{ a_{k}\ |\ k\ge n\}$.

      {\bf Solution}
\begin{enumerate}[(a)]
      \item {\bf To prove: }$\bigcap_{i}[a_{i}, b_{i}]\ne \emptyset$.
            \begin{proof}
                  Consider the sequences $(a_{n})$ and $(b_{n})$ independently.
                  Since $[a_{i}, b_{i}]$ always contains $[a_{i+1}, b_{i+1}]$, $(a_{n})$
                  must be monotone increasing, and $(b_{n})$ must be monotone decreasing.
                  Also, every term of $(b_{n})$ is an upper bound for $(a_{n})$, and
                  every term of $(a_{n})$ is a lower bound for $(b_{n})$. So,
                  $\alpha = \sup a_{n}$ and $\beta = \inf b_{n}$ must exist.\\

                  {\bf Claim: }$\alpha\leq \beta$.\\
                  FTSOC, assume $\beta<\alpha$. From the definition of supremum,
                  there must exist an $a_{i_{1}}$ such that $\beta<a_{i_{1}}<\alpha$.
                  From the definition of infimum, there must exist a $b_{i_{2}}$ such
                  that $\beta <b_{i_{2}}<a_{i_{1}}<\alpha$. This contradicts the fact that
                  every term of $(b_{n})$ is an upper bound for $(a_{n})$.\hfill $\Rightarrow\Leftarrow$\\

                  Thus, the interval $[\alpha, \beta]$ is always non empty.\\

                  {\bf Claim: }$\bigcap_{i}[a_{i}, b_{i}]= [\alpha, \beta]$\\
                  Every element in $[\alpha, \beta]$ is an upper bound for $(a_{i})$,
                  and a lower bound for $(b_{i})$. Thus, $[\alpha, \beta]\subset[a_{i}, b_{i}]$
                  for all $i$, i.e, $[\alpha, \beta]\subset \bigcap_{i}[a_{i}, b_{i}]$.\\
                  If $k\in \bigcap_{i}[a_{i}, b_{i}]$, $k$ is an upper bound of $(a_{n})$
                  and $k$ is a lower bound for $(b_{n})$. So, $k\ge \alpha$ and $k\le \beta$.
                  Thus, $k \in [\alpha, \beta]$.\\

                  Thus, $\bigcap_{i}[a_{i}, b_{i}]= [\alpha, \beta]\ne \emptyset$.
            \end{proof}

      \item {\bf To prove: }$|\bigcap_{i}[a_{i}, b_{i}]| =1\iff\lim_{ n \to \infty }(a_{n}-b_{n}) = 0$.
            \begin{proof}
                  We know that if a sequence is bounded and increasing,
                  it must converge to the least upper bound of its range.
                  Thus, $(a_{n})\to \alpha$. Similarly, if a sequence is bounded
                  and decreasing, it must converge to the greatest lower bound of its
                  range. Thus, $(b_{n})\to \beta$. Thus, we can say
                  \begin{align*}
                             & |\bigcap_{i}[a_{i}, b_{i}]| =1         \\
                        \iff & \alpha=\beta                           \\
                        \iff & \lim_{ n \to \infty }(a_{n}-b_{n}) = 0
                  \end{align*}
            \end{proof}
\end{enumerate}
\section{}
 {\bf Definitions}\\
$(a_{n})$ and $(b_{n})$ are sequences in $\mathbb{R}$.\\
$\sup a_{n} = \sup \{ a_{k}\ |\ k\ge n\}$.


      {\bf Solution}

\begin{lemma}$\sup(a_{n}+b_{n})\le \sup a_{n}+ \sup b_{n}$ for all $n\in \mathbb{N}$\end{lemma}
\begin{proof}
      If $\sup (a_n)$ or $\sup (b_n)$ is $\infty$, the inequality is trivially true.
      Thus, assume $\sup(a_n)$, $\sup(b_n) \in \mathbb{R}$.\\
      FTSOC, assume $\sup a_{n} +\sup b_{n}<\sup(a_{n}+b_{n})$.
      There must exist $a_{i}+ b_{i},\  i\ge n$ such that
      $\sup a_{n} +\sup b_{n}<a_{i}+b_{i}<\sup(a_{n}+b_{n})$. We know
      that $\sup a_{n}\ge a_{k}$ for all $k\ge n$. So, we have
      \begin{align*}
                     & (\sup a_{n}-a_{i}) +\sup b_{n}<b_{i} \\
            \implies & \sup b_{n}<b_{i},
      \end{align*}
      which is impossible.\hfill $\Rightarrow\Leftarrow$\\
\end{proof}

\begin{lemma}If $(a_{n})\to a\in \mathbb{R}\cup \{ \infty, -\infty \}$, $(b_{n})\to b\in \mathbb{R}\cup \{ \infty, -\infty \}$, and $a_{n}\le b_{n}$ for all $n\in \mathbb{N}$, then $a\le b$. \end{lemma}

\begin{proof}
      If $(b_{n})\to \infty$, the inequality is trivially true.
      If $(a_{n}) \to \infty$, it must be the case that $(b_{n})\to \infty$. Thus, assume
      $a, b\in \mathbb{R}$. \\
      FTSOC, assume $b<a$. Let $\epsilon = \frac{{a-b}}{2}$. Then, there exists an
      integer $N_{1}$ such that for all $n\ge N_{1}$, $|a_{n}-a|<\epsilon$. Similarly,
      there exists an integer $N_{2}$ such that for all $n\ge N_{2}$, $|b_{n}-b|<\epsilon$.
      Let $N$ be the greater of $N_{1}$ and $N_{2}$. Thus, for all
      $n\ge N$, $b_{n} \in \left( \frac{{3b-a}}{2}, \frac{{a+b}}{2} \right)$ and
      $a_{n}\in \left( \frac{{a+b}}{2}, \frac{{3a-b}}{2} \right)$, which implies
      $a_{n}> b_{n}$ for all $n\ge N$. \hfill $\Rightarrow\Leftarrow$\\
\end{proof}

{\bf To prove: }$\limsup_{ n \to \infty }(a_{n}+ b_{n})\le \limsup_{ n \to \infty}a_{n}+\limsup_{ n \to \infty } b_{n}$,
provided the sum on the right is not of the form $\infty-\infty$.

\begin{proof}
      Consider the sequences $(\sup a_{n}+\sup b_{n})$, and $(\sup(a_{n}+b_{n}))$. From Lemma 3.1,
      we have $\sup(a_{n}+b_{n})\le \sup a_{n}+ \sup b_{n}$ for all $n\in \mathbb{N}$.
      Also, it is clear the the suprema of tails of a sequence form a monotone decreasing sequence.
      Since all sequences are bounded in the extended reals, it follows these sequences must
      converge to a value in the extended reals. Thus, we can use Lemma 3.2:
      \begin{align*}
                     & \sup(a_{n}+b_{n})\le \sup a_{n}+ \sup b_{n}                                                                                                                          \\
            \implies & \lim_{ n \to \infty } (\sup(a_{n}+b_{n}))\le \lim_{ n \to \infty } (\sup a_{n}+ \sup b_{n}) = \lim_{ n \to \infty } (\sup a_{n})+ \lim_{ n \to \infty } (\sup b_{n}) \\
            \implies & \limsup_{ n \to \infty }(a_{n}+ b_{n})\le \limsup_{ n \to \infty}a_{n}+\limsup_{ n \to \infty } b_{n}
      \end{align*}
\end{proof}

\section{}
 {\bf Definitions}\\
$(p_{n})$ and $(q_{n})$ are Cauchy sequences in a metric space $(X, d)$.

      {\bf Solution}

      {\bf To prove: }The sequence $(d(p_{n}, q_{n}))$ converges.
\begin{proof}

      For any $m, n$, from the triangle inequality, we have
      \begin{align*}
            d(p_{n}, q_{n})\le d(p_{n}, p_{m}) + d(p_{m}, q_{n})
      \end{align*}
      and
      \begin{align*}
            d(p_{m}, q_{n})\le d(p_{m}, q_{m})+d(q_{m}, q_{n}).
      \end{align*}

      Combining the two, we get
      \begin{align*}
            d(p_{n}, q_{n})                            & \le d(p_{n}, p_{m}) + d(p_{m}, q_{m})+d(q_{m}, q_{n}) \\
            \implies |d(p_{n}, q_{n})-d(p_{m}, q_{m})| & \le d(p_{n}, p_{m})+d(q_{m}, q_{n}).
      \end{align*}
      Let $\epsilon>0$ be arbitrary. Since $(p_{n})$ and $(q_{n})$ are Cauchy sequences, there must exist integers $N_{1}$ and $N_{2}$ such that for all $m, n\geq N_{1}$, $d(p_{n}, p_{m})< \frac{\epsilon}{2}$ and for all $m, n\geq N_{2}$, $d(q_{n}, q_{m})< \frac{\epsilon}{2}$ respectively. Let $N$ be the larger of $N_{1}$ and $N_{2}$. Thus,
      \begin{align*}
            |d(p_{n}, q_{n})-d(p_{m}, q_{m})| < \frac{\epsilon}{2} + \frac{\epsilon}{2} = \epsilon
      \end{align*}
      for all $n\ge N$. This implies that the sequence $(d(p_{n}, q_{n}))$ is Cauchy. In $\mathbb{R}$, a sequence is Cauchy if and only if it is convergent (to a limit in $\mathbb{R}$).
\end{proof}

\section{}
 {\bf Definitions}\\
$(p_{n})$ is a sequence in $\mathbb{R}$. \\
For $a<b$, let $\zeta_{a, b}$ be the $[a, b]$-crossing number of $(p_{n})$.

      {\bf Solution}

      {\bf To prove: } $\zeta_{a, b}\in \mathbb{N}_{0} \ \ \forall \ a<b \iff (p_{n})\to p\in \mathbb{R}\cup \{ \infty, -\infty \}$.


\begin{proof}
      We will call a sequence with the property on the left (having finite $\zeta_{a,b}$ for all $a<b$) a \say{cross-bounded} sequence.

            {\bf Proof of convergent $\implies$ cross-bounded}

      We will prove the contrapositive: not cross-bounded $\implies$ not convergent.
      If a sequence is not cross-bounded, $\zeta_{a, b}=\infty$ for some $a<b$. This implies
      that for all $N$, there exists $d_{n}>N$ such that $p_{d_{n}}<a$,  and there exists $u_{n}>N$
      such that $p_{u_{n}}>b$. This allows us to make the following deductions:
      \begin{enumerate}[(i)]
            \item For $\epsilon=b-a$, there does not exist an $N$ such that for all $m, n\ge N$, $|p_{m}-p_{n}|<\epsilon$.
                  This tells us that $(p_{n})$ is not Cauchy, which implies $(p_{n})$ is not convergent in $\mathbb{R}$.
            \item For all $M>a$, there does not exist an $N$ such that for all $n\geq N$, $p_{n}\geq M$. Thus, $(p_{n})$ cannot converge to $\infty$.
            \item For all $m < b$, there does not exist an $N$ such that for all $n\geq N$, $p_{n}\leq m$. Thus, $(p_{n})$ cannot converge to $-\infty$.
      \end{enumerate}

      {\bf Proof of cross-bounded $\implies$ convergent}

      We can rephrase the notion of being cross-bounded like so: for every $a<b$, there
      either exists an $N$ such that for all $n\geq N$,  $p_{n}\ge a$, or there exists
      an $N$ such that for all $n\ge N$, $p_{n}\leq b$. In other words, eventually $a$
      becomes a lower bound for tails of $(p_n)$, or $b$ becomes an upper bound
      for tails of $(p_n)$. We will analyze bounded (in $\mathbb{R}$) and unbounded
      cases for $(p_{n})$ separately.
      \begin{enumerate}[(i)]
            \item Say $(p_{n})$ is bounded in $\mathbb{R}$, i.e, $m<p_{n}<M$ for all
                  $n\in \mathbb{N}$ for some $m, M\in \mathbb{R}$. Consider the interval
                  $$
                        (\alpha, \beta):=\left(\frac{{3m+M}}{4}, \frac{{m+3M}}{4}\right).
                  $$

                  Given our paraphrasing of what it means to be cross-bounded, there
                  either exists an $N$ such that for all $n\ge N$, $p_{n}\ge \alpha$, or
                  there exists an $N$ such that for all $n\ge N$, $p_{n} \leq \beta$. In both
                  cases, we have achieved new bounds for $(p_{n})$ for $n\ge N$, namely
                  $(\alpha, M)$ and $(m, \beta)$ respectively. Notice that
                  $M-a = \beta-m = \frac{3}{4}(M-m)$. Effectively, we have bounded a
                  tail of the sequence in an interval that is three-fourths the size of the
                  interval we started with. If we repeat this process $k$ times (always taking
                  the maximum of all the $N$'s we've chosen), we will bound a tail of
                  $(p_{n})$ in an interval of the size $\left( \frac{3}{4} \right)^{k}(M-m)$.
                  Taking the limit as $k\to \infty$,
                  $$
                        \lim_{ k \to \infty } (M-m)\left( \frac{3}{4} \right)^{k} = (M-m)\lim_{ k \to \infty }\left( \frac{3}{4} \right)^{k} = 0
                  $$
                  Thus, for arbitrary $\epsilon>0$, there will exist $k$ such
                  that $(M-m) \left( \frac{3}{4} \right)^{k}<\epsilon$, i.e, there will exist
                  a tail of $(p_{n})$ where the difference between any two terms is less
                  than $\epsilon$. This implies that $(p_{n})$ is Cauchy, and hence, convergent.
            \item If $(p_{n})$ is not bounded in $\mathbb{R}$, we will show that it converges
                  to $+\infty$ or $-\infty$. Consider any interval $(\alpha, \beta)$. As stated
                  previously, one of these should be true:
                  \begin{enumerate}[]
                        \item There exists an $N$ such that for all $n\ge N$, $p_{n} \geq \alpha$.

                              Since $(p_{n})$ is unbounded in $\mathbb{R}$ by hypothesis,
                              and $\alpha$ functions as a lower bound for all
                              $n\ge N$, $(p_{n})$ cannot have an upper bound. Thus, for any
                              interval $(\alpha', \beta')$, where $\alpha<\alpha'<\beta'$,
                              there must exist an $N'$ such that for all
                              $n\ge N'$,  $p_{n}\ge \alpha'$. Since $\alpha'$ was arbitrary,
                              this shows that $(p_{n})\to \infty$.
                        \item There exists an $N$ such that for all $n\ge N$, $p_{n} \leq \beta$.

                              Since $(p_{n})$ is unbounded in $\mathbb{R}$ by hypothesis,
                              and $\beta$ functions as an upper bound for all
                              $n\ge N$, $(p_{n})$ cannot have an lower bound. Thus, for any
                              interval $(\alpha', \beta')$, where $\alpha'<\beta'<\beta$,
                              there must exist an $N'$ such that for all $n\ge N'$,
                              $p_{n}\le \beta'$. Since $\beta'$ was arbitrary, this shows that 
                              $(p_{n})\to -\infty$.
                  \end{enumerate}

      \end{enumerate}


\end{proof}

\end{document}